% !TEX root = CoeneDylanBP.tex
%%=============================================================================
%% Intrinsieke Camerakalibratie
%%=============================================================================

\section{Intrinsieke camerakalibratie}
\label{sec:kalibratie}

Voordat een homografie kan worden berekend, moet de vervorming van de cameralens worden gecorrigeerd. Dit is een essentiële stap, aangezien een homografietransformatie veronderstelt dat rechte lijnen in de echte wereld ook als perfect rechte lijnen in het beeld verschijnen. In de praktijk introduceert vrijwel elke cameralens echter een zekere mate van optische vervorming, waardoor dit niet langer het geval is. Het doel van deze fase is dan ook om de \textit{intrinsieke cameraparameters} nauwkeurig te bepalen. Hiermee kan elk videoframe worden rechtgetrokken voordat de verdere beeldverwerking en positiebepaling plaatsvinden.

De gebruikte aanpak is gebaseerd op de methode van \textcite{Zhang2000}, die ook door OpenCV als standaard wordt gehanteerd. Hierbij wordt een schaakbordpatroon (\textit{checkerboard}) met bekende afmetingen gefilmd vanuit meerdere verschillende hoeken en afstanden. Omdat men exact weet hoe de geometrie van het schaakbord eruit hoort te zien, kunnen de afwijkingen in het camerabeeld worden gemeten. Het algoritme detecteert automatisch de hoekpunten van het schaakbord in elk frame en berekent op basis van de geobserveerde vervorming de volgende parameters:

\begin{itemize}
    \item De \textbf{camerakalibratiematrix} $K$, die de brandpuntsafstand ($f_x, f_y$) en het optische middelpunt ($c_x, c_y$) bevat:
          \[
              K = \begin{bmatrix}
                  f_x & 0   & c_x \\
                  0   & f_y & c_y \\
                  0   & 0   & 1
              \end{bmatrix}
          \]
    \item De \textbf{distortieco{\"e}ffici{\"e}nten}, die beschrijven hoe de lens het beeld vervormt.
\end{itemize}

Na de definitieve bepaling van deze parameters wordt elk frame van de video ontvormd met behulp van de functie \texttt{cv2.undistort} uit de OpenCV-bibliotheek. De kwaliteit van deze kalibratie wordt gemeten aan de hand van de \textit{reprojection error}. Dit is het gemiddelde verschil in pixels tussen de geobserveerde hoekpunten op de foto en de terug berekende positie van diezelfde punten volgens het model. Een waarde kleiner dan 1 pixel wordt hierbij als aanvaardbaar beschouwd voor nauwkeurige metingen.